\chapter{多源生态数据融合框架}

\section{数据来源与质量控制}

多源生态数据融合的首要任务是建立可追溯的数据登记表。示例将数据来源划分为现场调查、自动监测、遥感产品和管理资料四类。每一类数据都需要记录采集时间、空间范围、坐标系统、采样方法、字段含义和质量标记。表 \ref{tab:nwu-data-source} 展示了一个简化的数据来源清单。

\begin{table}[htbp]
  \centering
  \caption{多源生态数据来源示例}
  \label{tab:nwu-data-source}
  \begin{tabularx}{0.92\textwidth}{@{}lXX@{}}
    \toprule
    数据类型 & 典型字段 & 主要用途\\
    \midrule
    样线调查 & 物种名称、数量、样线编号、调查日期 & 提供物种出现和相对丰度信息\\
    红外相机 & 点位、触发时间、影像标签、独立事件 & 捕捉隐蔽动物活动节律\\
    遥感产品 & 植被指数、土地覆盖、海拔、坡度 & 描述栖息地背景条件\\
    管理资料 & 保护地边界、巡护记录、干扰事件 & 支撑保护措施解释\\
    \bottomrule
  \end{tabularx}
\end{table}

\section{特征融合流程}

特征融合应避免把来源差异误认为生态差异。示例流程包括坐标归一、时间窗口匹配、采样努力校正、变量标准化和不确定性记录五个步骤。图 \ref{fig:nwu-fusion-flow} 使用占位图展示模板对图题、交叉引用和正文段落的支持。

\begin{figure}[htbp]
  \centering
  \fbox{\parbox[c][4.2cm][c]{0.78\textwidth}{\centering 多源生态数据融合流程示意图}}
  \caption{多源生态数据融合流程}
  \label{fig:nwu-fusion-flow}
\end{figure}

对于正式论文，作者可以将占位图替换为真实流程图、研究区图或模型结构图。根据西北大学规范，图题置于图下方居中，表题置于表上方居中，图表编号采用全文连续编号。

\section{模型评价指标}

模型评价既要关注预测性能，也要关注解释稳定性。以栖息地适宜性预测为例，可以使用平均绝对误差、受试者工作特征曲线下面积和空间交叉验证误差等指标。若观测值为 $y_i$，预测值为 $\hat{y}_i$，则平均绝对误差定义为
\begin{equation}
  \operatorname{MAE} = \frac{1}{n}\sum_{i=1}^{n}\left|y_i-\hat{y}_i\right|.
\end{equation}
该公式用于验证西北大学模板中公式按章编号的排版效果。
